% !TEX program = xelatex
% 2025年度 (AY2025) 同志社大学大学院 生命医科学研究科 医工学コース
% 入学試験問題「数学」転記（問題文のみ。解答は 02_solutions/2025/ に置く）
% 原本: original/01_数学/2025/2025su.pdf
\documentclass[11pt,a4paper]{article}
\usepackage{../../../00_common/preamble}

\title{2025年度 (AY2025) 同志社大学大学院 生命医科学研究科\\
医工学コース 入学試験問題「数学」（転記）}
\author{}
\date{}

\begin{document}
\maketitle
\vspace{-3em}

\begin{rem}
本ファイルは原本スキャン \texttt{original/01\_数学/2025/2025su.pdf} の問題文を
転記したものである（解答は含めない）。判読は明瞭で,不明箇所はなかった。
\end{rem}

\bigskip

\noindent\textbf{〔I〕} 次の関数の $x=0$ 近傍で Taylor（テイラー）展開を求めよ.
\[
\sin x,\qquad -\infty<x<\infty
\]

\bigskip

\noindent\textbf{〔II〕} 次の極限値を求めよ.
\[
\lim_{n\to\infty}\iint_{D_n} e^{-x^2-y^2}\,dx\,dy,
\qquad
D_n=\bigl\{(x,y)\ ;\ x^2+y^2\le n^2,\ x\ge0,\ y\ge0\bigr\}
\]

\bigskip

\noindent\textbf{〔III〕} 次の行列の逆行列を求めよ.
\[
A=\begin{pmatrix}1&2&1\\0&1&0\\2&-1&1\end{pmatrix}
\]

\bigskip

\noindent\textbf{〔IV〕} 次の行列が対角化されるかを調べ,対角化できれば対角化せよ.
\[
B=\begin{pmatrix}4&1&0\\0&2&0\\0&0&3\end{pmatrix}
\]

\end{document}
